\chapter{Introduction}
% SMC background
% SQMC and its benefits
% what makes SQMC parallelizable 
% Football model and why SQMC works in this context
% Our contribution two parts - 1) GPU implementation of SQMC in cuthbert and 2) Modelling football matches with SQMC
% asynchronous, sparse local observations

Sequential Quasi-Monte Carlo (SQMC) is an extension of the Sequential Monte Carlo (SMC) algorithm, replacing the Monte Carlo proposal with deterministic propagation through low-discrepancy point sets. The error rate of SQMC is smaller than the Monte Carlo rate of $\mathcal{O}_P (N^{-1/2})$, and we can expect the algorithm to reach faster convergence for a smaller class of functions \citep{chopingerber2015sqmc}. Furthermore, recent work suggests that, for a fixed particle count, the error of a standard SMC exceeds any fixed threshold at infinitely many time instants with probability equal to one, while SQMC's error vanishes uniformly in time as $N$ grows \citep{gerber2026}. This safety guarantee, on a linear-Gaussian filtering problem, encourages the usage of SQMC over SMC in long-horizon problems.

However, the computational cost of SQMC is $\mathcal{O}(N \log N)$, where $N$ is the number of simulations at each iteration, a significant increase over SMC, which is $\mathcal{O}(N)$. This overhead stems from two additional steps in SQMC---sorting the particles along a Hilbert curve before resampling and applying the inverse-CDF transform to the quasi-random points. Another issue with SQMC is that its performance advantage over SMC decreases in increasing dimensions. The most likely factor is the curse of dimensionality that affects SMC \citep{chopingerber2017sqmcintroduction}. Importance sampling, a key component of the SMC algorithm, suffers from weight degeneracy in high dimensions: the larger the dimension, the greater the discrepancy between the proposal and the target distribution, so that a small fraction of particles carry non-negligible weight at each step \citep{snyderbengtssonbickel2008,bickellibengtsson2008}.

The dissertation makes two contributions. First, an implementation of the SQMC algorithm that is computationally efficient by leveraging accelerator hardware, specifically a GPU, using \texttt{jax}; the approach is generalisable to other similar parallel devices. We then show that this implementation outperforms current CPU implementations of the same algorithms. Second, we demonstrate how the high-dimensional limitations of SQMC can be overcome by exploiting the sparse observation structure of football matches. By modelling the transition distribution as linear-Gaussian, we can perform Rao-Blackwellisation to integrate out the latent states of non-playing teams, while performing SQMC on the smaller subset of playing teams, thereby reducing the 'effective dimension' of the model and mitigating the weight degeneracy problems as described above.

All code and experiments are available in the public repository \texttt{\seqsplit{ryantjx/sqmc-cuthbert}}; as part of an ongoing process of integrating the SQMC algorithm into the \texttt{cuthbert} project under issue \href{https://github.com/state-space-models/cuthbert/issues/255}{\#255}.